% SEGMENT OLP-0675-S01
% Part: proof-theory
% Chapter: normalization
% Section: permutations

\documentclass[../../../include/open-logic-section]{subfiles}

\begin{document}

\olfileid{pt}{nor}{per}

\olsection{வரிசைமாற்று உருமாற்றங்கள்}

\Log{N1i} இல் உள்ள வெட்டுகள், ஓர் !!a{elimination} விதியின்
முதன்முன்கோளில் முடியும் தொடர்கள். இயல்பாக்க நிறுவலின்
ஒரு நேர்வில், மிகுநிலை வெட்டுகளின் நீளத்தையும் அதன் மூலம்
!!{proof} இன் வெட்டு நீளத்தையும் குறைக்க வேண்டும்.
நீளம் $>1$ கொண்ட வெட்டு பல வழிகளில் முடிவடையலாம்.
அந்த வெட்டில் $!A$ இன் கடைசி !!{formula} நிகழ்வு
(\Elim{\lor} அல்லது \Elim{\lexists}) இன் முடிவா
என்பதையும், அது எந்த !!{elimination} அனுமானத்தின்
முதன்முன்கோளாக உள்ளது என்பதையும் இது பொறுத்தது.

% SEGMENT OLP-0675-S02
எடுத்துக்காட்டாக, தொடர்புடைய அனுமானங்கள் \Elim{\lor},
\Elim{\land} ஆகவும், $!A \ident !B \land !C$ ஆகவும்,
துணை-!!{proof}~$\delta_1$ பின்வருமாறு முடிவதாகவும் இருக்கட்டும்:
\[
\AxiomC{}
\RightLabel{$\delta_2$}
\DeduceC{$!D \lor !E$}
\AxiomC{$\Discharge{!D}{x}$}
\RightLabel{$\delta_3$}
\DeduceC{$!B \land !C$}
\AxiomC{$\Discharge{!E}{y}$}
\RightLabel{$\delta_4$}
\DeduceC{$!B \land !C$}
\DischargeRule{\Elim{\lor}}{x\,y}
\TrinaryInfC{$!B \land !C$}
\RightLabel{\Elim{\land}[1]}
\UnaryInfC{$!B$}
\DisplayProof
\]
வெட்டின் கடைசி !!{formula} நிகழ்வான~$!A_n$,
\Elim{\lor} இன் முடிவான கீழுள்ள $!B \land !C$ ஆகும்.
அதற்கு முந்தைய !!{formula} நிகழ்வான~$!A_{n-1}$,
\Elim{\lor} இன் சிறுமுன்கோள்களில் ஒன்றாகும்.

% SEGMENT OLP-0675-S03
இந்த வெட்டின் நீளத்தைக் குறைக்க, \Elim{\lor} மற்றும்
\Elim{\land} அனுமானங்களின் வரிசையை மாற்றலாம்
(வரிசைமாற்றலாம்). அதாவது $\delta$ இல் உள்ள
துணைநிறுவல்~$\delta_1$ ஐப் பின்வருவதாக மாற்றுகிறோம்:
\[
\AxiomC{}
\RightLabel{$\delta_2$}
\DeduceC{$!D \lor !E$}
\AxiomC{$\Discharge{!D}{x}$}
\RightLabel{$\delta_3$}
\DeduceC{$!B \land !C$}
\RightLabel{\Elim{\land}[1]}
\UnaryInfC{$!B$}
\AxiomC{$\Discharge{!E}{y}$}
\RightLabel{$\delta_4$}
\DeduceC{$!B \land !C$}
\RightLabel{\Elim{\land}[1]}
\UnaryInfC{$!B$}
\DischargeRule{\Elim{\lor}}{x\,y}
\TrinaryInfC{$!B$}
\DisplayProof
\]
\Elim{\lor} க்குக் கீழே உள்ள \Elim{\land} இன்
முன்கோளில் முடிந்த வெட்டுகள், இப்போது \Elim{\lor}
இன் சிறுமுன்கோள்களுக்கு மேலுள்ள \Elim{\land}
அனுமானங்களில் ஒன்றின் முன்கோளில் முடிகின்றன.
அதாவது ஒவ்வொன்றின் நீளமும்~$1$ குறைந்துள்ளது.

மேலும், $\delta_2$, $\delta_3$, $\delta_4$ இவற்றில்
முழுவதுமாகவோ, அல்லது $\delta$ இல் $\delta_1$ க்கு
முற்றிலும் வெளியிலோ அமைந்துள்ள எந்த வெட்டும் மாறாது:
அதே !!{formula}s இடம்பெறுகின்றன; அதே அனுமானங்கள்
வழியாகச் செல்கிறது. ஆகவே $\delta$ இன் ஒத்த வெட்டுக்கு
இருந்த அதே நிலையும் நீளமும் இதற்கும் உண்டு. எனவே
புதிய !!{proof} இன் வெட்டு நிலை உயரவில்லை;
புதிய !!{proof} இன் வெட்டு நீளம் பழைய
!!{proof} இன் வெட்டு நீளத்தைவிடக் குறைவு.

\begin{prob}
ஒரு வரிசைமாற்று உருமாற்றத்துக்குப் பின் குறைந்தது
ஒரு வெட்டின் நீளம்~$1$ குறைகிறது. \Elim{\land}
ஐ ஓர்~\Elim{\lor} வழியாக மேலே மாற்றினால்,
நீளம் குறையும் வெட்டுத் தொடர்கள் உண்மையில் இரண்டு;
ஆகவே இந்நேர்வில் !!{proof} இன் வெட்டு நீளம்
குறைந்தது~$2$ குறையும். ஒரே இத்தகைய வரிசைமாற்றால்
!!a{proof} இன் வெட்டு நீளம்~$2$ ஐவிட அதிகமாகக்
குறைய முடியுமா? வெட்டு \emph{நிலை} குறைய முடியுமா?
\end{prob}

ஓர் !!a{elimination} விதியை \Elim{\lor} அல்லது
\Elim{\lexists} விதிக்கு மேலே நகர்த்தும் செயல்
\emph{வரிசைமாற்று உருமாற்றம்} எனப்படும். சில விதி
இணைகளுக்கான வடிவங்கள்
\olref{tab:permutation-conversions} இல்
தரப்பட்டுள்ளன. (மீதியுள்ளவைப் பயிற்சிகளாக
விடப்பட்டுள்ளன.)

% SEGMENT OLP-0675-S04
% \begin{table}
% \begin{multline*}
% \shoveleft{\AxiomC{$!D \lor !E$}
% \AxiomC{$\Discharge{!D}{x}$}
% \shortDeduce
% \DeduceC{$!B \land !C$}
% \AxiomC{$\Discharge{!E}{y}$}
% \shortDeduce
% \DeduceC{$!B \land !C$}
% \DischargeRule{\Elim{\lor}}{x\,y}
% \TrinaryInfC{$!B \land !C$}
% \RightLabel{\Elim{\land}[1]}
% \UnaryInfC{$!B$}
% \DisplayProof}\\
% \shoveright{\longrightarrow\ 
% \AxiomC{$!D \lor !E$}
% \AxiomC{$\Discharge{!D}{x}$}
% \shortDeduce
% \DeduceC{$!B \land !C$}
% \RightLabel{\Elim{\land}[1]}
% \UnaryInfC{$!B$}
% \AxiomC{$\Discharge{!E}{y}$}
% \shortDeduce
% \DeduceC{$!B \land !C$}
% \RightLabel{\Elim{\land}[1]}
% \UnaryInfC{$!B$}
% \DischargeRule{\Elim{\lor}}{x\,y}
% \TrinaryInfC{$!B$}
% \DisplayProof}
% \\
% \shoveleft{\AxiomC{$!D \lor !E$}
% \AxiomC{$\Discharge{!D}{x}$}
% \shortDeduce
% \DeduceC{$!B \lif !C$}
% \AxiomC{$\Discharge{!E}{y}$}
% \shortDeduce
% \DeduceC{$!B \lif !C$}
% \DischargeRule{\Elim{\lor}}{x\,y}
% \TrinaryInfC{$!B \lif !C$}
% \AxiomC{$!B$}
% \RightLabel{\Elim{\lif}}
% \BinaryInfC{$!C$}
% \DisplayProof}\\
% \shoveright{\longrightarrow\ 
% \AxiomC{$!D \lor !E$}
% \AxiomC{$\Discharge{!D}{x}$}
% \shortDeduce
% \DeduceC{$!B \lif !C$}
% \AxiomC{$!B$}
% \RightLabel{\Elim{\lif}}
% \BinaryInfC{$!C$}
% \AxiomC{$\Discharge{!E}{y}$}
% \shortDeduce
% \DeduceC{$!B \lif !C$}
% \AxiomC{$!B$}
% \RightLabel{\Elim{\lif}}
% \BinaryInfC{$!C$}
% \DischargeRule{\Elim{\lor}}{x\,y}
% \TrinaryInfC{$!C$}
% \DisplayProof}
% \\
% \shoveleft{\AxiomC{$!D \lor !E$}
% \AxiomC{$\Discharge{!D}{x}$}
% \shortDeduce
% \DeduceC{$!B \lif !C$}
% \AxiomC{$\Discharge{!E}{y}$}
% \shortDeduce
% \DeduceC{$!B \lif !C$}
% \DischargeRule{\Elim{\lor}}{x\,y}
% \TrinaryInfC{$!B \lif !C$}
% \AxiomC{$!B$}
% \RightLabel{\Elim{\lif}}
% \BinaryInfC{$!C$}
% \DisplayProof}\\
% \shoveright{\longrightarrow\ 
% \AxiomC{$!D \lor !E$}
% \AxiomC{$\Discharge{!D}{x}$}
% \shortDeduce
% \DeduceC{$!B \lif !C$}
% \AxiomC{$!B$}
% \RightLabel{\Elim{\lif}}
% \BinaryInfC{$!C$}
% \AxiomC{$\Discharge{!E}{y}$}
% \shortDeduce
% \DeduceC{$!B \lif !C$}
% \AxiomC{$!B$}
% \RightLabel{\Elim{\lif}}
% \BinaryInfC{$!C$}
% \DischargeRule{\Elim{\lor}}{x\,y}
% \TrinaryInfC{$!C$}
% \DisplayProof}
% \\
% \shoveleft{\AxiomC{$!D \lor !E$}
% \AxiomC{$\Discharge{!D}{x}$}
% \shortDeduce
% \DeduceC{$\lforall[x][!B(x)]$}
% \AxiomC{$\Discharge{!E}{y}$}
% \shortDeduce
% \DeduceC{$\lforall[x][!B(x)]$}
% \DischargeRule{\Elim{\lor}}{x\,y}
% \TrinaryInfC{$\lforall[x][!B(x)]$}
% \RightLabel{\Elim{\lforall}}
% \UnaryInfC{$!B(t)$}
% \DisplayProof}\\
% \shoveright{\longrightarrow\ 
% \AxiomC{$!D \lor !E$}
% \AxiomC{$\Discharge{!D}{x}$}
% \shortDeduce
% \DeduceC{$\lforall[x][!B(x)]$}
% \RightLabel{\Elim{\lforall}}
% \UnaryInfC{$!B(t)$}
% \AxiomC{$\Discharge{!E}{y}$}
% \shortDeduce
% \DeduceC{$\lforall[x][!B(x)]$}
% \RightLabel{\Elim{\lforall}}
% \UnaryInfC{$!B(t)$}
% \DischargeRule{\Elim{\lor}}{x\,y}
% \TrinaryInfC{$!B(t)$}
% \DisplayProof}
% \end{multline*}
% \caption{Some permutation conversions for \Log{N1i}.}
% \ollabel{tab:permutation-conversions}
% \end{table}

{\small
\begin{longtable}{@{}l@{}}
\AxiomC{$!D \lor !E$}
\AxiomC{$\Discharge{!D}{x}$}
\shortDeduce
\DeduceC{$!B \land !C$}
\AxiomC{$\Discharge{!E}{y}$}
\shortDeduce
\DeduceC{$!B \land !C$}
\DischargeRule{\Elim{\lor}}{x\,y}
\TrinaryInfC{$!B \land !C$}
\RightLabel{\Elim{\land}[1]}
\UnaryInfC{$!B$}
\DisplayProof
$\longrightarrow$\\[6ex]
\quad\AxiomC{$!D \lor !E$}
\AxiomC{$\Discharge{!D}{x}$}
\shortDeduce
\DeduceC{$!B \land !C$}
\RightLabel{\Elim{\land}[1]}
\UnaryInfC{$!B$}
\AxiomC{$\Discharge{!E}{y}$}
\shortDeduce
\DeduceC{$!B \land !C$}
\RightLabel{\Elim{\land}[1]}
\UnaryInfC{$!B$}
\DischargeRule{\Elim{\lor}}{x\,y}
\TrinaryInfC{$!B$}
\DisplayProof
\\[10ex]
\AxiomC{$!D \lor !E$}
\AxiomC{$\Discharge{!D}{x}$}
\shortDeduce
\DeduceC{$!B \lif !C$}
\AxiomC{$\Discharge{!E}{y}$}
\shortDeduce
\DeduceC{$!B \lif !C$}
\DischargeRule{\Elim{\lor}}{x\,y}
\TrinaryInfC{$!B \lif !C$}
\AxiomC{$!B$}
\RightLabel{\Elim{\lif}}
\BinaryInfC{$!C$}
\DisplayProof
$\longrightarrow$\\[6ex]
\AxiomC{$!D \lor !E$}
\AxiomC{$\Discharge{!D}{x}$}
\shortDeduce
\DeduceC{$!B \lif !C$}
\AxiomC{$!B$}
\RightLabel{\Elim{\lif}}
\BinaryInfC{$!C$}
\AxiomC{$\Discharge{!E}{y}$}
\shortDeduce
\DeduceC{$!B \lif !C$}
\AxiomC{$!B$}
\RightLabel{\Elim{\lif}}
\BinaryInfC{$!C$}
\DischargeRule{\Elim{\lor}}{x\,y}
\TrinaryInfC{$!C$}
\DisplayProof
\\[10ex]
\AxiomC{$!D \lor !E$}
\AxiomC{$\Discharge{!D}{x}$}
\shortDeduce
\DeduceC{$!B \lif !C$}
\AxiomC{$\Discharge{!E}{y}$}
\shortDeduce
\DeduceC{$!B \lif !C$}
\DischargeRule{\Elim{\lor}}{x\,y}
\TrinaryInfC{$!B \lif !C$}
\AxiomC{$!B$}
\RightLabel{\Elim{\lif}}
\BinaryInfC{$!C$}
\DisplayProof
$\longrightarrow$\\[6ex]
\AxiomC{$!D \lor !E$}
\AxiomC{$\Discharge{!D}{x}$}
\shortDeduce
\DeduceC{$!B \lif !C$}
\AxiomC{$!B$}
\RightLabel{\Elim{\lif}}
\BinaryInfC{$!C$}
\AxiomC{$\Discharge{!E}{y}$}
\shortDeduce
\DeduceC{$!B \lif !C$}
\AxiomC{$!B$}
\RightLabel{\Elim{\lif}}
\BinaryInfC{$!C$}
\DischargeRule{\Elim{\lor}}{x\,y}
\TrinaryInfC{$!C$}
\DisplayProof
\\[10ex]
\AxiomC{$!D \lor !E$}
\AxiomC{$\Discharge{!D}{x}$}
\shortDeduce
\DeduceC{$\lforall[x][!B(x)]$}
\AxiomC{$\Discharge{!E}{y}$}
\shortDeduce
\DeduceC{$\lforall[x][!B(x)]$}
\DischargeRule{\Elim{\lor}}{x\,y}
\TrinaryInfC{$\lforall[x][!B(x)]$}
\RightLabel{\Elim{\lforall}}
\UnaryInfC{$!B(t)$}
\DisplayProof
\quad$\longrightarrow$\\[8ex]
\quad\AxiomC{$!D \lor !E$}
\AxiomC{$\Discharge{!D}{x}$}
\shortDeduce
\DeduceC{$\lforall[x][!B(x)]$}
\RightLabel{\Elim{\lforall}}
\UnaryInfC{$!B(t)$}
\AxiomC{$\Discharge{!E}{y}$}
\shortDeduce
\DeduceC{$\lforall[x][!B(x)]$}
\RightLabel{\Elim{\lforall}}
\UnaryInfC{$!B(t)$}
\DischargeRule{\Elim{\lor}}{x\,y}
\TrinaryInfC{$!B(t)$}
\DisplayProof
\\
\caption{\Log{N1i} இன் சில வரிசைமாற்று உருமாற்றங்கள்.}
\ollabel{tab:permutation-conversions}
\end{longtable}
}

% SEGMENT OLP-0675-S05
\begin{prob}
\Elim{\lor} ஐத் தொடர்ந்து \Elim{\lor} அல்லது
\Elim{\lnot} வரும்போது பொருந்தும் வரிசைமாற்று
உருமாற்றங்களைக் கொடுக்கவும்.
\end{prob}

\begin{prob}
\Elim{\lexists} ஐத் தொடர்ந்து \Elim{\exists}
வரும்போது பொருந்தும் வரிசைமாற்று உருமாற்றங்களைக்
கொடுக்கவும். பிந்தைய நேர்வில் எந்தத் தனிமாறி
நிபந்தனையும் ஏன் மீறப்படவில்லை என்பதை விளக்கவும்.
\end{prob}

% SEGMENT OLP-0675-S06
வரிசைமாற்று உருமாற்றத்தைப் பயன்படுத்தும்போது
அறியாமல் தனிமாறி நிபந்தனையை மீறிவிடவில்லை என்பதை
உறுதிசெய்துகொள்வோம். பின்வருவதைக் கருதுக:
\[
\AxiomC{}
\RightLabel{$\delta_2$}
\DeduceC{$!D \lor !E$}
\AxiomC{$\Discharge{!D}{x}$}
\RightLabel{$\delta_3$}
\DeduceC{$\lexists[x][!B(x)]$}
\AxiomC{$\Discharge{!E}{y}$}
\RightLabel{$\delta_4$}
\DeduceC{$\lexists[x][!B(x)]$}
\DischargeRule{\Elim{\lor}}{x\,y}
\TrinaryInfC{$\lexists[x][!B(x)]$}
\AxiomC{$\Discharge{!B(c)}{z}$}
\RightLabel{$\delta_5$}
\DeduceC{$!C$}
\DischargeRule{\Elim{\lexists}}{z}
\BinaryInfC{$!C$}
\DisplayProof
\]
% SEGMENT OLP-0675-S07
இதற்குப் பதிலாகப் பின்வருவதை வைத்தால்
\[
\AxiomC{}
\RightLabel{$\delta_2$}
\DeduceC{$!D \lor !E$}
\AxiomC{$\Discharge{!D}{x}$}
\RightLabel{$\delta_3$}
\DeduceC{$\lexists[x][!B(x)]$}
\AxiomC{$\Discharge{!B(c)}{z}$}
\RightLabel{$\delta_5$}
\DeduceC{$!C$}
\DischargeRule{\Elim{\lexists}}{z}
\BinaryInfC{$!C$}
\AxiomC{$\Discharge{!E}{y}$}
\RightLabel{$\delta_4$}
\DeduceC{$\lexists[x][!B(x)]$}
\AxiomC{$\Discharge{!B(c)}{z}$}
\RightLabel{$\delta_5$}
\DeduceC{$!C$}
\DischargeRule{\Elim{\lexists}}{z}
\BinaryInfC{$!C$}
\DischargeRule{\Elim{\lor}}{x\,y}
\TrinaryInfC{$!C$}
\DisplayProof
\]
% SEGMENT OLP-0675-S08
தனிமாறி~$c$, எடுத்துக்காட்டாக~$!D$ இல்
இடம்பெறலாம் என நாம் கவலைப்படலாம். ஏனெனில்,
முதலில் இருந்த~$\Elim{\lexists}$ க்கு மேலே
$!D$ எடுகோள் விடுவிக்கப்படுகிறது; அதனால் அசல்
!!{proof} இல் அந்த எடுகோள் \Elim{\exists}
அனுமானத்தின் முதன்முன்கோளுக்கு மேலே திறந்த
எடுகோளாக இல்லை. எனவே அங்கே தனிமாறி நிபந்தனை
நிறைவேறுகிறது. ஆனால் புதிய !!{proof} இல்,
\Elim{\lexists} அனுமானங்களுக்குப் பின்னரே $!D$
விடுவிக்கப்படுகிறது; ஆகவே அந்நிபந்தனை மீறப்படும்
போலத் தோன்றுகிறது. இருப்பினும் நாம் கருதும்
!!{proof}s ஒழுங்குறுவன என்பதை நினைவில் கொள்ளவும்:
ஒவ்வொரு தனிமாறியும் ஒரே ஓர் அனுமானத்தின் தனிமாறியாக
இருக்கிறது; முழு !!{proof} இலும் \Elim{\lexists}
இன் சிறுமுன்கோளுக்கு மேலே மட்டுமே இடம்பெறுகிறது.
குறிப்பாக, $c$ ஆனது $\delta_3$ இலோ $\delta_4$
இலோ இடம்பெற முடியாது.

வரிசைமாற்றத்தால் நகலாகும் துணைநிறுவலின் அனைத்து
தனிமாறிப் பெயர்களையும் ஒரு நகலில் புதிதாக மாற்றினால்தான்
புதிய நிறுவலும் ஒழுங்குறுவதாக இருக்கும்.

% SEGMENT OLP-0675-S09
இந்த எடுத்துக்காட்டு மேலும் ஒன்றைக் காட்டுகிறது:
புதிய !!{proof} இல்~$\delta_5$ இன் இரண்டு
நகல்கள் உள்ளன. எனவே $\delta_5$ இல் மிகுநிலை
வெட்டு இருந்தால், வெட்டு நீளத்தைக் குறைத்துவிட்டோம்
என்று \emph{சொல்ல முடியாது}! ஆகவே $\delta_5$
இல் மிகுநிலை வெட்டு இல்லை என்பது தெரிந்திருந்தால்தான்
வரிசைமாற்று உருமாற்றத்தைப் பயன்படுத்த வேண்டும்.
இதற்காக எப்போதும் \emph{மிக மேலுள்ள} மிகுநிலை
வெட்டைத் தேர்ந்தெடுப்போம்; அதாவது, அதன்
துணை-!!{proof} இல் (இந்நேர்வில்~$\delta_1$)
அந்தத் துணை-!!{proof} இன் கடைசி அனுமானத்துக்கு
மேலே முடியும் வேறெந்த மிகுநிலை வெட்டும்
இல்லாததாகும். இதனால்~$\delta_5$ இலும்
(உண்மையில் $\delta_2$, $\delta_3$, $\delta_4$
இவற்றிலும்) வேறு மிகுநிலை வெட்டுகள் இருக்கும்
வாய்ப்பு விலக்கப்படுகிறது.

\end{document}
